\section{Chellas neighborhood systems: complete source constructions}

\begin{annexitem}{annex:p8.cn.language}{thm:p08-cn-language}{Primitive possibility, definitional elimination and contextual replacement}
\begin{proof}
First derive diamond replacement using only classical principles.
If \(\vdash_L A\leftrightarrow B\), Boolean reasoning gives
\(\vdash_L\neg A\leftrightarrow\neg B\). RE yields
\(\Box\neg A\leftrightarrow\Box\neg B\); negate both
sides and use the two duality instances to obtain
\(\Diamond A\leftrightarrow\Diamond B\).
Replacement at a specified occurrence is then proved by induction
on the construction above that occurrence. At a Boolean parent,
equivalence of corresponding arguments implies equivalence of the
compound by a propositional tautology. At a box parent use RE;
at a diamond parent use the replacement just derived. Unchanged
arguments are equivalent to themselves. Several specified occurrences
are replaced one at a time; finitely many replacements compose by
transitivity of theorem equivalence. Zero replacements give reflexivity.
Every equivalence used in this argument is a theorem, not a conclusion
depending on local premises.

For the alternative classical characterization, assume the box
duality schema and diamond RE. Negating a theorem equivalence,
applying diamond RE and negating the result gives box RE by the
box duality schema. Taking its duality instance at \(\neg A\)
and replacing double negations under the diamond gives diamond
duality. Thus the original classical conditions hold. In the
forward direction the first argument already gives diamond RE,
and the original diamond duality at \(\neg A\), together
with box double-negation replacement, gives box duality.
Full contextual replacement implies either box or diamond RE
as its one-occurrence special case. Conversely the classical
conditions have already proved full contextual replacement.
This establishes both replacement-based characterizations.

Induct on \(A\) to prove \(A\leftrightarrow t(A)\).
Atoms and constants are fixed. Boolean cases follow by propositional
replacement of the induction equivalences. At a box, apply RE to its
inner induction equivalence. At a diamond, use
\(\Diamond B\leftrightarrow\neg\Box\neg B\), then
replace \(B\) by \(t(B)\) under negation and the box. This gives
exactly the diamond translation clause.

Translate a finite source theorem proof of E plus its chosen M/C/N
schemas. An RPL step is simulated by propositional conjunction
introduction for its finitely many inputs, the tautological implication
from their conjunction to its conclusion, and modus ponens. With
zero inputs it is just a tautology. The translated duality axiom
is a tautological self-equivalence. The
schemas M/C/N translate to the corresponding host instances. A source
RE step translates to a host RE step with its translated separate
theorem proof. Thus source theoremhood implies host theoremhood of
the translation. Conversely each host theorem proof, read in the
source's box-and-Boolean fragment, is already a source proof; its
rules and M/C/N instances occur in the source presentation.
For a local proof translate its premises and its imported theorems
and leave modus ponens unchanged. In the reverse direction use the
proved source equivalences between each original and translated
premise/conclusion. These finite Boolean steps recover the original
local consequence. Nothing in this proof permits modalizing a
merely premise-dependent equivalence.
\end{proof}
\end{annexitem}

\begin{annexitem}{annex:p8.cn.dual}{thm:p08-cn-dual}{Every Boolean dual case and the modality-word argument}
\begin{proof}
Induct on the formula to prove \(A^*\leftrightarrow\neg A\).
For an atom this is the definition; swapped constants give the
same equivalence. For negation, the induction equivalence gives
\((\neg B)^*=\neg B^*\leftrightarrow\neg\neg B\).
For conjunction, its dual is \(B^*\lor C^*\), equivalent to
\(\neg B\lor\neg C=\neg(B\land C)\); disjunction is
the other De Morgan case. The implication dual
\(\neg B^*\land C^*\) becomes \(B\land\neg C\),
the negation of \(B\to C\). The equivalence dual
\(B^*\leftrightarrow\neg C^*\) becomes
\(\neg B\leftrightarrow C\), propositionally equivalent
to \(\neg(B\leftrightarrow C)\).
For a box the dual is \(\Diamond B^*\), equivalent by
contextual replacement to \(\Diamond\neg B\), and by
duality and double-negation replacement to \(\neg\Box B\).
For a diamond use \(\Box B^*\leftrightarrow\Box\neg B
\leftrightarrow\neg\Diamond B\).
All replacements here are theorem equivalences from the induction.
Applying the result to \(A^*\) gives
\(A^{**}\leftrightarrow\neg A^*\leftrightarrow A\).

For words, induct on their length. The empty word gives
\(A\leftrightarrow\neg\neg A\). Suppose
\(\nu A\leftrightarrow\neg\nu^*\neg A\).
Prepending a box and applying box replacement gives
\(\Box\nu A\leftrightarrow\Box\neg\nu^*\neg A
\leftrightarrow\neg\Diamond\nu^*\neg A\).
Prepending a diamond is the symmetric duality argument; prepending
negation follows by Boolean double-negation cancellation.
This proves \(\mu A\leftrightarrow\neg\mu^*\neg A\).

If \(\mu A\to\nu A\) is a theorem for every \(A\),
these identities and contraposition give
\(\nu^*\neg A\to\mu^*\neg A\). Set \(A=\neg B\)
and replace double negations to get the implication for every \(B\).
The converse uses the same argument and \(\mu^{**}=\mu\)
for words. Applying it to both implication directions proves
the corresponding equivalence assertion.
\end{proof}
\end{annexitem}

\begin{annexitem}{annex:p8.cn.semantics}{thm:p08-cn-semantics}{Complete truth-set and atom-agreement inductions}
\begin{proof}
Induct on formulas. Atomic truth sets are the specified valuations,
truth is \(W\), and falsity is empty. Classical pointwise negation
gives complement, conjunction gives intersection and disjunction union.
An implication is false just when its antecedent is true and consequent
false, giving \((W\setminus\|A\|)\cup\|B\|\); equivalence
is the intersection of the two implication sets.
The neighborhood box clause gives truth set \(f(\|A\|)\).
Its diamond clause gives \(W\setminus f(W\setminus\|A\|)\).
These are exactly the operation presentation's recursive clauses,
so they agree at every induction stage and world.

The diamond truth clause is the clause for \(\neg\Box\neg A\),
so the source duality is valid. If \(A\leftrightarrow B\)
holds at all worlds of a model, their truth sets are equal. Membership
of those sets in each \(N(w)\) is therefore equal as well.
This proves global boxed equivalence in that model, and also
preservation of validity across any class of models. It does not
follow from equality only at a single world.

For atom agreement, compare two valuations on the same \(W,N\).
At an occurring atom their truth sets agree by hypothesis. Constants
agree unconditionally; each Boolean truth-set operation preserves
equalities established at its arguments. At a box or diamond, the
same operation on equal inner truth sets gives equal output sets.
This proves agreement at every world for the whole formula.
Empty models have no worlds at which a formula can fail and no
pointed witnesses, so they add no restriction to validity or
pointed local consequence.
\end{proof}
\end{annexitem}

\begin{annexitem}{annex:p8.cn.closure}{thm:p08-cn-closure}{Closure commutation and exact frame validation}
\begin{proof}
The intersection-decomposition condition implies supplementation:
if \(X\in\mathcal B\) and \(X\subseteq Y\), then
\(X=X\cap Y\in\mathcal B\) forces \(Y\in\mathcal B\).
Conversely a family closed upwards contains both supersets \(X,Y\)
of any member \(X\cap Y\). Thus the conditions are equivalent.
Binary intersection closure yields nonempty finite intersection
closure by induction on the number of members; the converse is
its two-member case.

If \(Z\in(\mathcal B^+)^-\), write
\(Z=\bigcap_{i=1}^nY_i\), \(n\geq1\), with each
\(Y_i\supseteq X_i\in\mathcal B\). Then
\(Z\supseteq\bigcap_iX_i\), so \(Z\in(\mathcal B^-)^+\).
Conversely, if \(Z\supseteq\bigcap_iX_i\), with
\(X_i\in\mathcal B\) and \(n\geq1\), each
\(Z\cup X_i\in\mathcal B^+\) and distributivity gives
\(\bigcap_i(Z\cup X_i)=Z\cup\bigcap_iX_i=Z\).
Hence \(Z\in(\mathcal B^+)^-\). Empty input families
give empty families on both sides, as only positive arities are used.

If \(\mathcal B\) is intersection-closed, two members of
\(\mathcal B^+\) contain respective \(X,Y\in\mathcal B\).
Their intersection contains \(X\cap Y\in\mathcal B\),
so it belongs to \(\mathcal B^+\). If \(\mathcal B\) is
supplemented and \(Z\) contains an intersection of its members
\(X_i\), each \(Z\cup X_i\) already belongs to \(\mathcal B\).
The same distributivity identity shows \(Z\in\mathcal B^-\).
This proves the two preservation assertions. A nonempty supplemented
family contains \(W\), since it contains a subset of \(W\);
a family containing \(W\) is nonempty.

For the schema correspondence, put \(X=\|A\|\), \(Y=\|B\|\).
Supplementation makes \(X\cap Y\in N(w)\) imply both
memberships, validating M. Binary intersections validate C by
the reverse implication. Unit membership validates N.
For necessity of these frame conditions, arbitrary valuations can
assign any \(X,Y\) to two distinct atoms. A missing intersection
closure then directly falsifies C. A missing upward inclusion has
\(X\in N(w)\), \(X\subseteq Y\) and \(Y\notin N(w)\).
Assign \(p\) to \(Y\) and \(q\) to \(X\): their conjunction
has set \(X\), so M is false at \(w\). A missing unit falsifies
\(\Box\topv\).

The valuation qualification matters. On three worlds let every
family contain only \(\{a\}\), and give all atoms the empty
valuation. Every formula truth set is then either empty or all worlds:
the Boolean clauses preserve these two values, boxes are false
everywhere and diamonds true everywhere. M and C are valid under
this valuation, although the actual family is not supplemented.
Finally, on one world \(\mathcal B=\{\varnothing\}\),
supplementation changes \(\Box\topv\) from false to true.
On two worlds with \(\mathcal B=\{\{a\},\{b\}\}\),
intersection closure changes the box of the conjunction of their
two singleton-valued atoms from false to true. Quasi-filtering
does so too. These show why closure operations cannot be claimed
to preserve arbitrary formula truth before a preservation proof.
\end{proof}
\end{annexitem}

\begin{annexitem}{annex:p8.cn.rules}{thm:p08-cn-rules}{Every theorem-rule direction, the zero case and iteration}
\begin{proof}
If a classical logic has RM, applying it to both directions of a
theorem equivalence gives RE. Applying RM to the two conjunction
projections and conjoining the resulting implications gives M.
Conversely if M and RE hold, a theorem \(A\to B\) gives
\(A\leftrightarrow(A\land B)\) by propositional reasoning.
RE identifies their boxes. M implies
\(\Box(A\land B)\to\Box B\), so \(\Box A\to\Box B\).
Thus RM and M are equivalent in classical systems.

RR gives RM by applying it to the theorem
\((A\land A)\to B\) whenever \(A\to B\) is a theorem,
then contracting the identical boxed antecedents propositionally.
RR applied to the reflexive conjunction implication gives C.
Conversely RM and C turn a theorem \((A\land B)\to C\)
into \(\Box(A\land B)\to\Box C\), then use C to infer
\((\Box A\land\Box B)\to\Box C\). This proves
regularity exactly M+C and the names \(M=EM,R=EMC\).

A regular logic has the K schema: C combines
\(\Box(A\to B)\) and \(\Box A\) into the box of
\((A\to B)\land A\); RM on the tautology from that conjunction
to \(B\) gives \(\Box B\). In a monotonic logic with the
K schema, RM on \(A\to(B\to A\land B)\) gives
\(\Box A\to\Box(B\to A\land B)\), and K combines
that consequent with \(\Box B\) to give C. Thus C and
the K schema are equivalent in monotonic systems, without N.

In a regular logic, induction gives boxed conjunction introduction
for any positive number of conjuncts. RM then carries a theorem
\(\bigwedge_{i=1}^nA_i\to B\) to the boxed implication.
For zero conjuncts that theorem is just \(B\); if N holds,
\(B\leftrightarrow\topv\) is a theorem and RE together
with N yields \(\Box B\). This is exactly the missing
zero-antecedent part of RK. It also proves N implies theorem-only
necessitation, while necessitation of truth gives N in the reverse
direction.

For the usual K-schema/necessitation calculus, necessitate a theorem
implication and apply K to obtain RM. This obtains RE and M by
the first argument. Necessitate the tautology
\(A\to(B\to A\land B)\) and apply K twice to obtain C.
Necessitating truth gives N. Conversely M/C/N/RE already provide
K and necessitation, so the theorem presentations are identical.
Their local consequences, which add only premises and modus ponens,
are therefore identical as well.

Diamond RM follows by contraposing the input implication, applying
box RM, contraposing again and using duality. For diamond RR,
a theorem \(A\to B\lor C\) gives
\((\neg B\land\neg C)\to\neg A\). Box RR and
contraposition give \(\Diamond A\to\Diamond B\lor\Diamond C\).
Its converse follows by the same transformation. The diamond
meet/join schemas follow either from these rules on tautologies
or from the duality lemma applied to their box counterparts.

For depth \(k\), repeated RE gives equivalence replacement under
\(\Box^k\), and repeated RM gives implication replacement.
Applying this to conjunction projections gives its M schema at
every depth. Regularity gives binary boxed-meet equivalence;
if depth \(k\) has that equivalence, one application of RE
and the outer binary-meet equivalence gives depth \(k+1\).
Thus regularity gives C, RR and K at every finite depth as well.
Depth zero is the propositional case. Cancellation of consecutive
complements proves \(\Diamond^k A\leftrightarrow
\neg\Box^k\neg A\) by induction, and duality gives the
corresponding iterated diamond assertions. None of these steps
modalizes a premise-dependent implication.
\end{proof}
\end{annexitem}

\begin{annexitem}{annex:p8.cn.separation}{thm:p08-cn-separation}{Soundness and the complete eight-system separation}
\begin{proof}
Each \(L_S\) is sound on its required frame class. Induct on
finite theorem proofs: propositional tautologies hold pointwise,
modus ponens preserves truth, the selected schemas are valid by
the closure theorem, and RE preserves validity because its separate
theorem equivalence is valid at every world. The induction may
equivalently use the total size of a proof tree including each RE
subproof. The one-world principal family \(\{W\}\) validates
all three schemas and has false falsity, so the eight systems are
consistent.

The first witness's empty family is upward and intersection-closed
vacuously, but omits \(W\), falsifying N. In the second,
the two singleton sets and \(W\) are all nonempty subsets of
the two-world domain. Their family is upward closed and contains
the unit. The two singleton-valued atoms are both boxed, while
their conjunction has empty extension, omitted from the family.
Thus C is false. In the third, \(\{\varnothing,W\}\)
is intersection-closed and unit-containing, but the conjunction
of the singleton-valued atoms is boxed while neither atom is,
so M is false.

Adding schemas gives \(L_S\subseteq L_T\) when
\(S\subseteq T\). Otherwise choose a schema in \(S\setminus T\).
Its witness above satisfies both other conditions, hence all
conditions demanded by \(T\). Soundness makes that schema a
non-theorem of \(L_T\), while it is a theorem of \(L_S\).
This proves the reverse inclusion criterion and separates every
pair of the eight systems, including incomparable pairs.

For the third model, \(\Box\bot\) and
\(\Box(\bot\to p)\) are true because their sets are empty
and \(W\), while \(\Box p\) is false for a singleton.
The K schema is therefore false in this EN-model. The empty-family
witness separates EMC from adding N. It makes every \(\Box A\)
false irrespective of \(A\), and is a model for each of E, EM,
EC and EMC; soundness excludes all such boxed theorems in them.
In the second model, empty is absent, so
\(\Diamond\topv\), equivalent to \(\neg\Box\bot\),
is true everywhere. A singleton and its complement are both
members, so \(\Box p\land\Box\neg p\) is true.
This verifies the deontic distinction without importing intersection
closure from a stronger system.
\end{proof}
\end{annexitem}

\begin{annexitem}{annex:p8.cn.maximal}{thm:p08-cn-maximal}{Every premise, extension and proof-set case}
\begin{proof}
For deduction, transform a proof from \(\Gamma\cup\{B\}\)
into one of \(B\to A\) from \(\Gamma\). A theorem or
\(\Gamma\)-premise \(X\) gives \(B\to X\) by the
tautology \(X\to(B\to X)\). The extra premise \(B\)
gives \(B\to B\). A modus-ponens step follows from the
two transformed implications using the propositional distribution
tautology. A modal replacement line has an independent theorem
proof and is therefore a theorem line in this argument. The
reverse follows by adding \(B\) and applying modus ponens.
Every proof uses finitely many premises. Repeated deduction thus
identifies local derivability with theoremhood of a finite-premise
conjunction implication.

For the source's original extension construction enumerate all
formulas \(A_1,A_2,\ldots\), allowing duplicates, and put
\(\Delta_0=\Gamma\). At stage n add \(A_n\) if
\(\Delta_{n-1}\cup\{A_n\}\) is consistent; otherwise
leave \(\Delta_n=\Delta_{n-1}\). Every stage is consistent
and the stages increase. Their union \(\Delta\) extends
\(\Gamma\). Every finite subset of that union is in a single
stage, by taking a maximum of the finitely many stages at which
its elements first appear. A contradiction proof from the union
would thus contradict consistency of that stage.
If \(\Delta\cup\{B\}\) is consistent, choose n with
\(A_n=B\). Its subset \(\Delta_{n-1}\cup\{B\}\)
is consistent, so that stage adds B and B is in the union.
This proves maximality of the original skip-on-rejection construction.
The source's additional stage lemma also holds: if \(A_k\)
is in the union, it is already in \(\Delta_k\). Otherwise
it was rejected there because \(\Delta_{k-1}\cup\{A_k\}\)
was inconsistent, and no later consistent stage could contain
that set. Repeated occurrences in the enumeration do not affect
this argument.

One may instead add \(\neg A_n\) when A is rejected, as in
the source's alternative exercise. Deduction gives that negation
from the preceding set; its addition cannot destroy consistency,
for that would also give the positive formula from the preceding
set and contradict its consistency. The same finite-union argument
then proves maximality. Both are countable existence constructions,
not effective consistency tests for an arbitrary classical system.

A maximal set is deductively closed. If a deducible A were absent,
maximality would make adjoining A inconsistent; deduction would
give the negation of A from the same set. The two deductions
would make the original set inconsistent, a contradiction.
It therefore contains theorems and is closed under modus ponens.
If A is absent, that same inconsistent-extension argument gives
its negation, which belongs by deductive closure. Both A and its
negation cannot belong to a consistent set. Thus every formula
is decided exactly once, and negation is membership complement.
Truth belongs to the set and falsity does not.
For conjunction use the two projections and propositional
conjunction introduction. For disjunction, either disjunct entails
it; if neither belongs, both negations belong and contradict
a supposed disjunction. An implication belongs exactly when
its antecedent is absent or its consequent present: these two
cases imply it propositionally, and an implication with present
antecedent yields the consequent by modus ponens. Equivalence
belongs exactly when the two component memberships agree, using
its two implications and the two propositional equal-truth cases.
These prove all Boolean proof-set identities.

A theorem implication gives proof-set inclusion by modus ponens
in every maximal set. If \(A\to B\) is not a theorem,
\(\{A,\neg B\}\) is consistent, for otherwise deduction
and classical reasoning would give that implication. Extend it
to a maximal set; this world belongs to \(|A|_L\) and not
\(|B|_L\), disproving inclusion. Equality is therefore equivalent
to both theorem implications, hence to theorem equivalence.
RE gives boxed equivalence for equal proof sets, and maximal-set
closure makes their box memberships equal in every world.
\end{proof}
\end{annexitem}

\begin{annexitem}{annex:p8.cn.canonical}{thm:p08-cn-canonical}{All canonical and closure-preserving truth cases}
\begin{proof}
Equality of proof sets implies theorem equivalence, hence boxed
equivalence by RE. Thus the defining rule for membership of a
proof set cannot depend on which formula represents it. A
canonical family exists: take the smallest family, and optionally
add non-proof sets. Those additions cannot change a proof-set
membership.

Induct on formulas for the truth lemma. Atoms have their specified
proof sets. Truth/falsity and all Boolean cases are exactly the
maximal membership conditions just proved. At a box, induction
identifies the inner truth set with its proof set, whose neighborhood
membership is equivalent by definition to box membership in the
world. At a diamond, the complemented inner proof set is
\(|\neg A|_L\). Its absence from the family is equivalent
to \(\Box\neg A\notin u\), hence to
\(\neg\Box\neg A\in u\), and by duality to
\(\Diamond A\in u\). This proves the whole truth lemma.

If C is a theorem schema, two smallest-family members are
\(|A|_L,|B|_L\) with \(\Box A,\Box B\in u\).
C and propositional closure give \(\Box(A\land B)\in u\).
Its proof set is their intersection, so the smallest family is
intersection-closed. If N is a theorem, \(\Box\topv\in u\)
and \(|\topv|_L=W_L\), so every canonical family contains
the unit.

If M holds, supplement the smallest family. When a proof set
\(|A|_L\) enters this supplementation, it contains some
\(|B|_L\) with \(\Box B\in u\). The inclusion gives
\(\vdash_L B\to A\), and RM gives
\(\vdash_L\Box B\to\Box A\). Thus \(\Box A\in u\)
already. Conversely its box membership already supplied its
smallest-family membership. Supplementation therefore preserves
every proof-set requirement and remains canonical. It supplies
M, retains N when required, and retains C when required by the
closure theorem. Choosing the smallest or supplemented model
according to the presence of M gives all eight stated constructions,
with their full truth lemmas.
\end{proof}
\end{annexitem}

\begin{annexitem}{annex:p8.cn.completeness}{thm:p08-cn-completeness}{Strong local completeness and the invalid local RE witness}
\begin{proof}
Theorem soundness was proved in the separation annex. In a local
proof, imported theorems are true at every world of every required
frame; its premises are true at the chosen world and modus ponens
preserves that truth. Thus local derivability implies local
semantic consequence under precisely the stated premise convention.

If \(\Gamma\nvdash_{L_S}A\), the set
\(\Gamma\cup\{\neg A\}\) is consistent: deduction
would otherwise give \(\Gamma\vdash_{L_S}A\).
Extend it to a maximal world \(u\). The appropriate canonical
construction has all required M/C/N properties and makes exactly
the formulas of \(u\) true there. It satisfies \(\Gamma\)
and falsifies \(A\), so semantic consequence fails. This gives
the reverse direction for any premise set, and theoremhood
determination at empty premises. If consequence holds, the
resulting syntactic proof uses a finite subset of its premises;
soundness gives semantic consequence from that finite subset.
Equivalently, finitely satisfiable premise sets are consistent and
are satisfied in the same maximal-model construction. This proves
the compactness assertions.

For the local-RE boundary take worlds \(a,b\), with
\(R(a)=R(b)=\{b\}\), and use their principal neighborhoods.
This is a K-frame. Put \(V(p)=\{a\}\) and \(V(q)=W\).
At \(a\) the two atoms are both true, so their equivalence is
true. At its sole successor \(b\), \(p\) is false and \(q\)
true, so \(\Box p\) is false and \(\Box q\) true at \(a\).
Their boxed equivalence fails. Modal replacement from that local
premise, or necessitation of a local truth, would therefore be
unsound even on this normal frame.
\end{proof}
\end{annexitem}

\begin{annexitem}{annex:p8.cn.augmentation}{thm:p08-cn-augmentation}{Principal representation and the source normal canonical argument}
\begin{proof}
For \(N_R(w)=\{X:R(w)\subseteq X\}\), the intersection
of its members is exactly \(R(w)\): the set \(R(w)\)
is itself a member and is contained in every other member.
So the frame is augmented. Conversely an augmented frame is
exactly the principal frame for successors \(\bigcap N(w)\).
With the same atom valuation, induction on formulas gives equal
truth in these paired models. Boolean cases preserve equal
inner truth values. The box case is the equivalence of neighborhood
membership with inclusion of the successor set in the inner truth
set. The diamond case is the complement of that test, equivalently
existence of a successor in the inner truth set, including the
empty-successor case. This proves pointwise equivalence.

On a finite world domain each unit-containing neighborhood family
is finite and nonempty. Binary intersection closure therefore
makes the intersection of all its members another member.
Supplementation gives all supersets of that core. Every finite
filter is thus augmented, including improper filters whose core
is empty and whose family is all subsets.

For the source infinite example, at real \(a\) take all sets
containing some \((a,b)\) with \(b>a\). The family is upward
closed and contains \(\mathbb R\). Two members contain intervals
\((a,b_1),(a,b_2)\); their intersection contains
\((a,\min(b_1,b_2))\), proving intersection closure.
Every real \(x\leq a\) is outside each such interval; for
\(x>a\), choose \(a<b<x\) to exclude it from one member.
Thus the core is empty. No member is empty, since every included
interval has a midpoint. Principal families include their core,
so none can represent this family on the same worlds for all
valuations. An atom can name any subset and detect a differing
neighborhood membership.

For canonical augmentation in a normal \(L\), let
\(B_u=\{A:\Box A\in u\}\). The core of its smallest
canonical family is exactly \(\{v:B_u\subseteq v\}\).
If \(\Box A\in u\), every core world contains \(A\),
so augmentation contains \(|A|_L\). If \(\Box A\notin u\),
\(B_u\cup\{\neg A\}\) is consistent. Otherwise a finite
subset \(B_1,\ldots,B_n\) would give the theorem
\(\bigwedge_iB_i\to A\). The proved normal rule RK,
including \(n=0\), would give the boxed conjunction implication;
the boxes in \(u\) would yield \(\Box A\in u\), a contradiction.
Extend that consistent set to a maximal \(v\). It is in the core
and outside \(|A|_L\). Augmentation therefore adds no incorrect
proof-set membership, and is canonical. For K it supplies every
needed countermodel while principal frames validate M/C/N, proving
its determination both by filters and by augmented frames. The
infinite example shows why agreement of these logics is a different
claim from identity of their unrestricted frame classes.
\end{proof}
\end{annexitem}

\begin{annexitem}{annex:p8.cn.interval}{thm:p08-cn-interval}{The actual interval-membership distinction}
\begin{proof}
With only exact open intervals around \(a\) in its family, the
box clause is literally the assertion that the entire formula
truth set is such an interval. Supplementation instead tests
whether some such interval is a subset of that truth set.
At zero assign an atom the set \((-1,1)\cup\{3\}\).
It contains \((-1,1)\), an open interval around zero, so the
supplemented box is true. It is not an interval: it contains
zero and three but not two. The original box is therefore false.
Existence of an open interval around a point on which the atom
is true is exactly the latter, containment-based local-persistence
condition. It is not equivalent to requiring that the complete
truth set itself be an interval.
\end{proof}
\end{annexitem}

\begin{annexitem}{annex:p8.cn.filtration}{thm:p08-cn-filtration}{All primitive modal clauses and every closure case}
\begin{proof}
For each \(A\in F\), its original truth set is a union of
equivalence classes by definition of the quotient. Such a saturated
set \(X\) satisfies \(\pi^{-1}[\pi[X]]=X\).
Complements and finite intersections of saturated sets are saturated,
and their images are the respective complements and intersections
of images. In particular this applies to all observed propositions,
including complements arising from primitive diamonds.

Worlds in the same quotient class agree on each observed box truth,
and on each observed diamond truth. The latter is the negation of
membership of the observed complement. Hence membership of every
observed proposition in the original family agrees at those worlds,
so the finest filtered family is well defined. If two observed
propositions have the same image, taking preimages shows they were
the same set, so they cannot give inconsistent membership prescriptions.
For an observed \(X\), its image belongs to the finest family
exactly when \(X\in N(w)\): a generating equal image has the
same saturated preimage.

Induct on \(A\in F\). Atoms have their saturated images;
constants have the whole quotient or empty image. Complement,
intersection, union and implication images match the Boolean truth
clauses. For \(\Box A\), induction identifies its inner filtered
truth set with \(\pi[\|A\|]\), and the established membership
equivalence gives the same box truth. For \(\Diamond A\),
the inner complement is \(\pi[W\setminus\|A\|]\);
its nonmembership gives the same diamond truth. This proves the
finest filtration truth lemma in every case.

We next verify that each additional closure retains the prescribed
membership for every observed \(X\). Left to right holds because
all closures include their generating families. For the reverse
in the supplemented case, membership of \(\pi[X]\) supplies
a generator \(\pi[Y]\subseteq\pi[X]\) with \(Y\in N(w)\).
Preimages give \(Y\subseteq X\); original supplementation
then gives \(X\in N(w)\).
For intersection closure, \(\pi[X]\) is an intersection of
finitely many generating images \(\pi[Y_i]\), with positive
arity. Preimages give \(X=\bigcap_iY_i\), which belongs to
the original intersection-closed family. For quasi-filtering,
closure commutation writes membership as
\(\bigcap_i\pi[Y_i]\subseteq\pi[X]\), again with positive
arity. Preimages give \(\bigcap_iY_i\subseteq X\).
Original finite intersection closure and supplementation give
\(X\in N(w)\). Thus no closure produces an incorrect
observed membership or nonmembership, and the same induction
proves its whole \(F\)-truth lemma.

The closures have their stated structural properties by definition
and the closure theorem. If the original family contains the unit,
the observation \(\Box\topv\in F\) inserts \(\pi[W]=W_F\)
in each finest family. All closures retain it, preserving N.
Finally a class is specified by the truth pattern on \(F\).
Distinct classes have distinct patterns, so their number is at
most \(2^{|F|}\). Empty original models cause no pointed
failure; filtering a pointed countermodel yields at least one class.
\end{proof}
\end{annexitem}

\begin{annexitem}{annex:p8.cn.decision}{thm:p08-cn-decision}{Finite determination and both terminating effective searches}
\begin{proof}
If \(A\) is not a theorem, the completeness construction gives
a pointed countermodel in its required M/C/N class. Choose the
finite subformula set in the statement, adding the unit observation
where needed. The relevant filtered/closed construction has all
required properties and still falsifies \(A\), on at most
\(2^{|F|}\) worlds. Soundness supplies the other direction of
finite determination.

Compute this integer bound from the input syntax. For each
\(1\leq n\leq2^{|F|}\), label the worlds
\(0,\ldots,n-1\). Enumerate all subsets of the world set,
all families of those subsets at every world, and all valuations
of the finitely many input atoms. These are finite lists. Check
M by all subset/superset pairs, C by all intersections of pairs
of members, and N by the full set's membership, imposing only
the selected conditions. Evaluate the formula bottom up by its
Boolean and modal truth-set clauses. Every step is a finite
membership or Boolean-set operation.

If a retained table has a falsifying world, return that table,
valuation and world; soundness certifies a non-theorem. If all
tables up to the bound pass, a non-theorem is impossible by
the bounded countermodel result. Thus the procedure halts with
the correct answer. Any finite countermodel can be renamed to
one of these labeled domains; induction on formulas shows a
bijection preserving valuations and neighborhoods preserves truth.
Unmentioned atoms can be empty by the atom-agreement lemma.

For the source's dovetail route, enumerate finite theorem proof
objects by size. Propositional tautology checks are effective
finite truth-table checks with modal constituents treated as
propositional arguments. Schema instances, modus ponens and
RE with its finite independent theorem subproof are effectively
checkable. This lists every theorem of the finitely presented
systems. In parallel enumerate finite models by size and their
finite tables, as above. A theorem eventually has a proof and
cannot have a sound countermodel; a non-theorem eventually has
a finite countermodel. Alternating these two searches therefore
terminates in exactly one direction. For finite premises, the
deduction theorem reduces local consequence to theoremhood of
their conjunction implication. This does not give a decision
algorithm for an arbitrarily supplied infinite set of premises.
\end{proof}
\end{annexitem}

\begin{annexitem}{annex:p8.cn.algebra}{thm:p08-cn-algebra}{Boolean soundness, the logical quotient and every modal condition}
\begin{proof}
For a powerset algebra, choose \(\ast=f\) and the original
atom sets. The truth-set induction already proved identifies
algebraic values with formula truth sets, and the Boolean unit
is \(W\). All-world truth is therefore exactly unit value.

Propositional tautologies evaluate to the unit in any Boolean
algebra, without a representation assumption. To see this directly,
take the finitely many Boolean values \(b_1,\ldots,b_m\)
occurring as propositional arguments. For each bit pattern make
the minterm consisting of the meet of \(b_i\) when its bit
is one and \(\neg b_i\) when zero. Distributivity makes
these minterms join to one; different ones meet to zero.
On a minterm a Boolean formula is either that minterm or zero,
according to its ordinary bit-pattern truth value, by induction
on the Boolean operations. A tautology therefore includes every
minterm and evaluates to one. Modus ponens preserves unit value
because \(1\to b=b\). If an equivalence has unit value,
its two values are equal; applying the function \(\ast\)
to equal values preserves equality, proving RE soundness. The
defined diamond uses \(\neg\ast\neg\) and satisfies duality.
The three algebraic conditions make the corresponding M/C/N
schemas unit-valued. This proves theorem soundness for each system.

The relation \(\vdash_L A\leftrightarrow B\) is an
equivalence relation by propositional reflexivity, symmetry and
transitivity. Boolean operations on its classes are well defined
because propositional equivalence is a congruence. Every Boolean
algebra identity is represented by a propositional theorem, so
the quotient is a Boolean algebra with unit \([\topv]_L\)
and zero \([\bot]_L\). RE makes the modal operation on
classes well defined too. Induction gives canonical value
\([A]_L\) for every formula, including primitive possibility:
its value is \([\neg\Box\neg A]_L=[\Diamond A]_L\)
by duality. A class equals the unit exactly when
\(A\leftrightarrow\topv\) is a theorem, which by
propositional reasoning is exactly theoremhood of \(A\).

The Boolean order on classes satisfies
\([A]_L\leq[B]_L\) exactly when \(A\to B\) is a
theorem. For each required M/C schema, its instances therefore
give the displayed inequality on arbitrary classes, every one
of which is represented by a formula. N gives \(\ast1=1\).
Thus a non-theorem has a non-unit canonical algebraic value
in an algebra with all required conditions. Together with the
soundness argument this proves all eight determination results.
It does not equate a quotient class with a set of possible worlds.
\end{proof}
\end{annexitem}

\begin{annexitem}{annex:p8.cn.algebra-finite}{thm:p08-cn-algebra-finite}{Finite projection, all three preserved conditions and the necessary restriction}
\begin{proof}
The Boolean algebra generated by \(m\) elements is finite.
The minterms in the preceding proof have at most \(2^m\)
nonzero members, and every element of the generated algebra
is a join of a subset of them. Hence there are at most
\(2^{2^m}\) such elements. Applied to the finitely many
subformula values, this proves finiteness of \(B_F\).

The join defining \(r(b)\) exists because its candidate set
is finite and contains zero. It belongs to \(B_F\), lies
below \(b\), and is its greatest \(B_F\)-lower bound.
In particular \(r(b)=b\) for \(b\in B_F\), and
\(r(1)=1\). The map is monotone. It preserves binary meet:
\(r(b)\land r(c)\) is a \(B_F\)-element below
\(b\land c\), so lies below \(r(b\land c)\);
the latter lies below both \(r(b)\) and \(r(c)\),
giving the reverse inequality.

For \(a,b\in B_F\), if the original operation satisfies M,
monotonicity and meet preservation of \(r\) give
\[
 r(\ast(a\land b))\leq r(\ast a\land\ast b)
                   =r(\ast a)\land r(\ast b).
\]
If it satisfies C, the reversed original inequality gives the
reversed projected inequality. If it satisfies N, then
\(\ast_F1=r(\ast1)=r(1)=1\). Each selected combination
is therefore retained by the finite operation.

Assign occurring atoms their original values, which are in
\(B_F\), and assign other atoms arbitrarily in that algebra.
Induct on subformulas for preservation. Boolean operations are
those of the subalgebra, so preserve the induction values.
If \(\Box B\) is a subformula, both \(v(B)\) and
\(\ast v(B)=v(\Box B)\) are represented in \(B_F\).
Thus \(\ast_Fv(B)=r(v(\Box B))=v(\Box B)\).
If \(\Diamond B\) is a subformula, its original value
\(\neg\ast\neg v(B)\) is in \(B_F\); so is its
complement \(\ast\neg v(B)\). The projection fixes that
complement and retains the diamond value. Constants are retained
because the subalgebra has the same unit and zero. The root's
non-unit value remains non-unit, giving the finite countermodel.

Agreement on every subalgebra argument is stronger and can be
impossible. Take \(B=\mathcal P(\{a,b,c\})\) and a
subformula set consisting only of an atom with value \(\{a\}\).
Its Boolean subalgebra is
\(\{\varnothing,\{a\},\{b,c\},W\}\).
An allowed arbitrary modal operation can send \(\{a\}\)
to \(\{b\}\), which is outside that subalgebra. No operation
from the subalgebra into itself can agree at that argument.
The projection instead retains every modal value actually required
by the subformula induction; it makes no impossible agreement
claim about unrepresented outputs.
\end{proof}
\end{annexitem}

\begin{annexitem}{annex:p8.cn.other-conditions}{thm:p08-cn-other-conditions}{All five schemas, their converses and general iteration}
\begin{proof}
Write \(X=\|A\|\). D requires that membership of \(X\)
imply nonmembership of its complement, exactly (d). T requires
that membership of \(X\) at \(w\) imply \(w\in X\),
exactly (t). B requires \(w\in X\) to imply membership
of the diamond extension \(d(X)\), exactly (b). The extension
of \(\Box A\) is \(f(X)\), so 4 requires (iv).
For 5, \(\Diamond A\) at \(w\) means
\(W\setminus X\notin N(w)\); (v) applied to that
complement makes \(d(X)\in N(w)\), which is
\(\Box\Diamond A\). Conversely, if the schema 5 is
valid for all valuations, give an atom extension
\(W\setminus Y\) for any \(Y\notin N(w)\).
It forces \(W\setminus f(Y)\in N(w)\), proving (v).
The other four converses likewise assign an atom any desired
set \(X\); their pointwise clauses then yield the frame
condition. This use of arbitrary valuations is explicit.

By induction, the box iterate has truth set \(f^n(X)\).
The diamond operation is complement, then \(f\), then
complement. Its \(n\)-fold iterate is complement, then
\(f^n\), then complement, because the two middle complements
cancel at each composition; the zeroth case is the identity.
Substituting these sets into the antecedent and consequent of
the displayed G-schema gives exactly the inclusion in the
statement. A formula implication holds at all worlds precisely
when its antecedent set is included in its consequent set.
Arbitrary atom valuations give necessity of the all-\(X\)
condition, and the truth-set calculation gives sufficiency.
The converse implication reverses the inclusion. Diamond dual
schema correspondences also follow from the proved syntactic
duality, without imposing relational conditions on an arbitrary
neighborhood frame.
\end{proof}
\end{annexitem}

\begin{annexitem}{annex:p8.cn.t5-principles}{thm:p08-cn-t5-principles}{All ET5 deductions, six-way separation and its non-normal counterexample}
\begin{proof}
T at falsity gives \(\Box\bot\to\bot\), hence
\(\neg\Box\bot\), equivalent to \(P=\Diamond\topv\).
Schema 5 gives \(P\to\Box P\), so \(\Box P\) is
a theorem. Since \(P\leftrightarrow\topv\) is then a
theorem, RE gives \(\Box P\leftrightarrow\Box\topv\)
and hence N. The earlier theorem-rule proof supplies RN.
For any theorem \(A\), diamond replacement of
\(A\leftrightarrow\topv\), together with P, gives
\(\Diamond A\), the source's RP.

Contraposing T at \(\neg A\) gives \(A\to\Diamond A\).
Combine it with 5 to obtain B, \(A\to\Box\Diamond A\).
Combine it with T to obtain D. T applied to \(\Diamond A\)
and schema 5 give \(\Box\Diamond A\leftrightarrow\Diamond A\).
Its modality dual is \(\Diamond\Box A\leftrightarrow\Box A\).
Now B at \(\Box A\) gives
\(\Box A\to\Box\Diamond\Box A\).
Replace its inner \(\Diamond\Box A\) by \(\Box A\)
using the proved theorem equivalence; this yields 4.
Together with T it gives \(\Box\Box A\leftrightarrow\Box A\).
The dual gives \(\Diamond\Diamond A\leftrightarrow\Diamond A\).

Any positive operator word of length at least two can shorten
its adjacent pair using these four equivalences and theorem
replacement. It reduces to box or diamond. Move any internal
negation outward using \(\Box\neg A\leftrightarrow
\neg\Diamond A\) and \(\Diamond\neg A\leftrightarrow
\neg\Box A\), then cancel paired negations. Thus every
modality reduces to identity, box or diamond, optionally negated.
For exact distinctness take two worlds with the universal
accessibility relation, hence principal family \(\{W\}\)
at both. It satisfies (t) and (v). Give an atom extension
\(\{a\}\). At \(a\), box is false and identity true;
at \(b\), identity is false and diamond true. Box and
diamond differ at both. Negating these evaluations separates
the negative three. Setting the atom everywhere true separates
every positive modality from every negative one. Hence ET5
has exactly six. T and its dual supply the indicated forward
implications, while the same witnesses refute their reverses.

For non-normality let \(W=\{a,b,c\}\) and
\[
 \mathcal C=\{\varnothing,W,\{a\},\{b,c\},\{b\},\{a,c\}\}.
\]
It is complement-closed. Define \(f(X)=X\) for
\(X\in\mathcal C\), and \(f(X)=\varnothing\) otherwise,
with \(N(w)=\{X:w\in f(X)\}\). Always \(f(X)\subseteq X\),
so (t) holds. Also \(f(W\setminus f(X))=W\setminus f(X)\):
if X is in the family its complement is too; otherwise
\(f(X)\) is empty and \(f(W)=W\). This gives (v).
Assign \(p\) extension \(\{a,b\}\) and \(q\) extension
\(\{a,c\}\). At \(a\), \(\Box(p\land q)\) is true
because the intersection \(\{a\}\) is in \(\mathcal C\),
but \(\Box p\) is false. M fails. ET5 is sound on this
frame, and every normal system contains M, so ET5 is not normal.
The modal reduction count alone does not identify it with S5.
\end{proof}
\end{annexitem}

\begin{annexitem}{annex:p8.cn.t5-completeness}{thm:p08-cn-t5-completeness}{Each of the three actual source canonical constructions}
\begin{proof}
Soundness follows by induction on theorem proof objects, using
the corresponding T/5 validation from the preceding condition
proof and the extensional RE argument. The one-world principal
identity frame validates both, so the systems are consistent.

For ET, any member \(X\) of its smallest canonical family
is \(|A|_L\) with \(\Box A\in u\). T implies
\(A\in u\), hence \(u\in X\). This proves (t).
For E5 use the largest canonical family. If \(X\notin N(u)\),
X must be a proof set \(|A|_L\), with \(\Box A\notin u\).
At every world \(v\), membership of that same proof set is
equivalent to \(\Box A\in v\). Thus
\[
                 \{v:X\notin N(v)\}=|\neg\Box A|_L.
\]
By duality, schema 5 yields theorem negative introspection
\(\neg\Box A\to\Box\neg\Box A\). Its antecedent
belongs to \(u\); the consequent belongs too, putting the
displayed negative extension in the canonical family. This is (v).

For ET5 the smallest family satisfies (t) as above. If an
omitted X is a proof set, the same negative-introspection
argument gives its negative extension as a forced proof-set
member. If X is not a proof set, it is absent from every
smallest family, so its negative extension is \(W_L\).
The derived N theorem puts \(W_L=|\topv|_L\) in every
family. This establishes (v) also for all non-proof sets.

Each is therefore a canonical model in the specified frame
class. Its full truth lemma and the maximal-set construction
give a countermodel for every failed theorem. They give a
point satisfying \(\Gamma\) and \(\neg A\) whenever
local derivability fails, exactly as in the earlier strong
completeness proof. Soundness gives the reverse directions.
No claim about a generic smallest canonical model for every
other combination of additional schemas is used.
\end{proof}
\end{annexitem}

\begin{annexitem}{annex:p8.cn.nonproof}{thm:p08-cn-nonproof}{Omitted-set necessity and an explicit smallest-E5 failure}
\begin{proof}
In any canonical model of an E5-extension, an omitted proof
set \(X=|A|_L\) has negative extension
\(|\neg\Box A|_L\), since its membership at each world
is exactly that world's box membership. Negative introspection
puts the box of that formula at every world where X is absent,
so its proof set belongs to the corresponding neighborhood
family. Condition (v) is consequently automatic on proof sets.
For non-proof sets it asserts precisely the omitted-set
restriction in the statement. It imposes no obligation from
the mere inclusion of such a set. The largest family includes
all non-proof sets, so its remaining obligations are empty.

Here is a complete construction showing why the smallest
E5 choice is different. On worlds \(a,b\), put
\(N(a)=\{\{a\}\}\) and \(N(b)=\mathcal P(W)\).
At b no subset is absent, so (v) is vacuous. If X is absent
at a, it differs from \(\{a\}\); X is present at b and
absent at a, so \(\{v:X\notin N(v)\}=\{a\}\in N(a)\).
Thus this frame satisfies (v), and E5 is sound on it.
At a the unit W is absent, so \(\Box\topv\) is false.
The set \(u^\circ\) of all formulas true at that point is
E5-maximal: it contains all E5 theorems, is modus-ponens closed,
decides every formula, and cannot derive falsity from its true
premises using the fixed local proof convention. It contains
\(\neg\Box\topv\). Hence the smallest canonical family
at this world omits \(W_L\).

To construct a non-proof set without a cardinality assumption,
enumerate all formulas \(A_n\). For each n take the one-world
identity frame with exactly atom \(p_n\) true and the other
atoms false. This is an E5-frame. Let \(u_n\) be its complete
true theory, an E5-maximal set by the same argument. These worlds
are pairwise distinct because of their atom valuations. Put
\[
                       X=\{u_n:A_n\notin u_n\}.
\]
For every n, membership of \(u_n\) in X is the opposite of
its membership in \(|A_n|_L\). Thus X differs from every
proof set, since every formula occurs in the enumeration.
In the smallest canonical model X belongs to no neighborhood
family at any world. Its negative extension is therefore
\(W_L\). At \(u^\circ\) that extension is not in the
family, proving a failure of (v).

This example satisfies any vacuous restriction placed only on
the included extra non-proof sets, because there are none.
The successful E5-largest and ET5-smallest constructions above
instead meet the required restriction on omitted sets. The
counterexample leaves their completeness proofs intact and
supplies the exact missing distinction.
\end{proof}
\end{annexitem}
